% !TEX root = ../main.tex
\section{Introduction}
\label{sec:intro}

When a language model outputs the sentence ``The Eiffel Tower is in Paris,''
what has happened, epistemically speaking? The question sounds naive, and the
answers pull in incompatible directions in a way that matters for anyone who
has to build, evaluate, or sign off on one of these systems. On one answer, now
standard in both technical and popular discourse, the model knows---or at least
believes---that the Eiffel Tower is in Paris; its output expresses that
attitude; and when its outputs go wrong the model is mistaken, or lying, or
\emph{hallucinating}, in the term of art that has colonised the field. On
another answer, favoured by deflationary critics, the model believes nothing
whatever: it is a stochastic text generator, and epistemic vocabulary applied to
it is metaphor at best.

Both answers are underdescribed, because neither engages the formal structure of
the attitudes being ascribed or withheld. The received tradition in formal
epistemology does not treat \emph{knowledge} and \emph{belief} as labels to be
applied wherever behaviour looks apt. It treats them as operators with
characteristic logical profiles. Knowledge, on the standard idealisation, is
factive and normal; belief is normal, consistent, and introspective. Both
validate the rule of congruentiality (RE), which requires that logically
equivalent contents be treated alike. That requirement is not decoration. It
encodes what it is, structurally, for a state to be a belief-like attitude at
all: that the attitude is directed at a \emph{proposition}---an entity
individuated no more finely than logical equivalence---and that it stands in at
least idealised rational relations to the system's other attitudes.

My first thesis is that model assertion, regimented as an operator $\Out_M$,
fails these conditions, and fails them constitutively rather than marginally.
$\Out_M$ fails K, M, D, 4, and 5. Critically, it fails RE: a model may assert
$\phi$ under one formulation and deny it under a logically equivalent
formulation, as a matter of architecture rather than noise. RE-failure is the
signature of hyperintensionality, and a hyperintensional operator admits no
Kripke semantics---no adjustment of frame conditions helps, because every Kripke
frame validates RE. Nor, as I show in Section~\ref{sec:hyperintensionality},
does bare neighborhood semantics suffice, since neighborhoods are sets of truth
sets and logically equivalent formulas share a truth set. What remains is a
hyperintensionalised neighborhood model or an impossible-worlds semantics. Since
knowledge and belief in the received sense are normal, and $\Out_M$ is not even
classical, language models are neither epistemic nor doxastic systems in that
sense. The best available characterisation is that they are \emph{sub-doxastic}
systems in a sense descended from \citet{stich1978}---though the descent
involves a mutation I state rather than elide.

It is worth being clear about which half of that conclusion is hard-won. That
model output is not \emph{knowledge} follows from factivity alone: models assert
falsehoods, so $\Out_M \phi \rightarrow \phi$ is invalid, and no modal
sophistication is required to see it. The apparatus earns its keep on the
doxastic claim, which is the contested one.

My second thesis concerns hallucination. The term is standardly defined
truth-conditionally: a hallucination is a false or fabricated output. This gets
the phenomenon wrong in both directions. Define a grounding operator
$\Grd_M(\phi \mid c)$: $\phi$ is supported by the model's evidential base,
relative to a practical context $c$ fixing jurisdiction, time, and intended use.
Then hallucination is ungrounded assertion, and error is grounded false
assertion, with no truth condition appearing in the first definition. An
assertion faithfully transmitted from an erroneous source is an error, not a
hallucination. An ungrounded assertion that happens to be true is a
hallucination regardless---a failure of safety, in the sense familiar from
\citet{sosa1999}, masked by luck. The 2$\times$2 taxonomy that results shows
that truth-matching benchmarks misclassify both off-diagonal cells, and
Proposition~\ref{prop:independence} shows this is not a matter of degree:
accuracy places no bound whatever on the grounding rate.

The context index on $\Grd_M$ is not decoration either, and it is where the
philosophical apparatus meets the problem as it actually presents itself in
regulated work. A claim grounded in a source that was authoritative in 2019 and
superseded in 2024 is not thereby grounded for a decision taken now. A claim
grounded in European regulation is not grounded for a decision taken under
another jurisdiction. Without the index, both cases collapse into ``error'' and
get charged to a defective evidential base---when the base was not defective but
inapplicable, and the remedy is currency and applicability checking rather than
corpus correction.

\paragraph{Plan.} Section~\ref{sec:preliminaries} compresses the formal
preliminaries. Section~\ref{sec:output-operator} defines $\Out_M$ and
demonstrates the axiom failures. Section~\ref{sec:hyperintensionality} develops
the hyperintensionality result and its semantic consequences.
Section~\ref{sec:contextual-truth} introduces contextual truth and the failure
mode I call illicit universalization, which motivates the context index.
Section~\ref{sec:grounding} defines $\Grd_M$ and the taxonomy.
Section~\ref{sec:hallucination} develops hallucination as unwarranted assertion
and proves the independence result. Section~\ref{sec:world-access} addresses
what training exposure does and does not supply.
Section~\ref{sec:subdoxastic} argues for the sub-doxastic characterisation and
states the limits of the analogy. Section~\ref{sec:internal} confronts the
literature on internal truth representations, which bears directly on the
argument's central premise. Section~\ref{sec:objections} considers objections.
Section~\ref{sec:evaluation} draws out implications for evaluation, and
Section~\ref{sec:warrant} sketches the transition to institutional warrant.

\paragraph{Scope.} Two notes. First, my claims concern current autoregressive
transformer models as deployed, including instruction-tuned, preference-optimised,
and outcome-reinforcement-trained reasoning variants. The last of these deserve
particular attention, because they are the systems most plausibly moving toward
inferential closure, and I take up the question they raise in
Section~\ref{sec:reasoning}. Nothing here rules out future architectures whose
assertion operators are normal, or nearly so; the framework says precisely what
such an architecture would have to achieve, and
Section~\ref{sec:evaluation} states it as a checklist. Second, the argument is
not that these systems are epistemically worthless. Sub-doxastic systems can be
extraordinarily useful, and I build them for a living. The argument is that
their usefulness must be theorised, and their failures measured, under the
correct formal description---and that the description currently in use produces
metrics that reward the wrong behaviour.
